% Work/MAT – 1. ročník – hodina 50. Navazuje na 49 (vytýkání a seskupování).
% Primární metodický zdroj: Realisticky.cz, 010808 Rozklad mnohočlenů na součin III,
% uživatelem dodaný dokument „08 Rozklad mnohočlenů na součin III.pdf“, zejména Př. 2–8.
% Doplňující zdroj: Work-MAT/ucebnice/01_Zakladni_poznatky/Algebraicke_vyrazy/Algebraické výrazy.tex.
% Úlohy jsou nově formulované a výběrově upravené podle témat a postupů zdroje.
\documentclass[11pt,a4paper]{article}
% Spolecny styl pro ucitelske pripravy MAT (Work/MAT)
% Verze 1.1, 2026-10-23
\NeedsTeXFormat{LaTeX2e}
\ProvidesPackage{priprava}[2026/10/23 v1.1 Sablona ucitelskych priprav MAT]

% Zakladni typografie a cestina
\RequirePackage[T1]{fontenc}
\RequirePackage[utf8]{inputenc}
\RequirePackage[czech]{babel}
\RequirePackage{lmodern}
\RequirePackage{microtype}
\RequirePackage{amsmath,amssymb}
\RequirePackage{xcolor}
\RequirePackage{enumitem}
\RequirePackage[most]{tcolorbox}
\RequirePackage{tabularx,booktabs,array}
\RequirePackage{fancyhdr}
\RequirePackage{needspace}
\RequirePackage{geometry}
\RequirePackage[hidelinks]{hyperref}

\geometry{a4paper,top=17mm,bottom=18mm,left=23mm,right=23mm}
\definecolor{notegray}{gray}{0.39}
\setlength{\parindent}{0pt}
\setlength{\parskip}{2.5pt plus 1pt minus 0.5pt}
\setlist[itemize]{topsep=2pt,itemsep=1pt,parsep=0pt,leftmargin=1.45em}
\setlist[enumerate]{topsep=2pt,itemsep=2pt,parsep=0pt,leftmargin=1.65em}
\setcounter{secnumdepth}{0}
\renewcommand{\arraystretch}{1.12}

% Zahlavi: vlevo rocnik | globalni tema, vpravo cislo: nazev hodiny
% Zapati zustava zcela prazdne.
\pagestyle{fancy}
\fancyhf{}
\renewcommand{\headrulewidth}{0.35pt}
\renewcommand{\footrulewidth}{0pt}
\setlength{\headheight}{14pt}
\newcommand{\nastavPripravu}[4]{%
  \fancyhead[L]{\small #1\enspace\textbar\enspace #2}%
  \fancyhead[R]{\small #3: #4}%
}

% Minimalisticke cernobile boxy. Kazda definice zvlast.
\tcbset{commonlesson/.style={enhanced,breakable=false,colback=white,colframe=black,
  colbacktitle=white,coltitle=black,boxrule=0.45pt,
  arc=0.8mm,outer arc=0.8mm,left=3mm,right=3mm,top=1.6mm,bottom=1.6mm,
  toptitle=1mm,bottomtitle=0.7mm,fonttitle=\bfseries,
  attach boxed title to top left={xshift=2.5mm,yshift=-1.15mm},
  boxed title style={colback=white,colframe=white,boxrule=0pt,left=1mm,right=1mm}}}
\newtcolorbox{zadaniBox}[1][]{commonlesson,title={Zadání},#1}
\newtcolorbox{prikladBox}[1][]{commonlesson,title={Příklad},#1}
\newtcolorbox{definiceBox}[1][]{commonlesson,title={Definice},#1}

% Spolecny box pro zadani a jeho reseni, bez mezery mezi castmi.
% Pouziti:
%   \begin{prikladReseniBox}
%     Zadani ulohy.
%     \tcblower
%     \begin{align*} ... \end{align*}
%   \end{prikladReseniBox}
% Spodni cast ma jemny sedy podklad a delici linku; lze v ni sazet
% vice radku matematiky. Obe casti maji jen jeden spolecny ramecek.
\tcbset{solutionlesson/.style={
  commonlesson,
  enhanced jigsaw,
  breakable,
  colbacklower=black!3,
  segmentation style={solid,draw=black!25,line width=0.35pt},
  fontlower=\small\color{notegray},
  before lower={\noindent\textbf{Řešení:}\enspace},
  leftlower=3mm,rightlower=3mm
}}
\newtcolorbox{prikladReseniBox}[1][]{solutionlesson,title={Příklad},#1}
\newtcolorbox{zadaniReseniBox}[1][]{solutionlesson,title={Zadání},#1}

% Pomucky pro autora pripravy: sedive poznamky; \reseni pro starsi pripravy
\newcommand{\poznamka}[1]{\par\begingroup\small\color{notegray}\itshape
  \textnormal{\textbf{Poznámka pro učitele.}}\enspace #1\par\endgroup}
\newcommand{\reseni}[1]{\par\begingroup\small\color{notegray}
  \textbf{Očekávané řešení:}\enspace #1\par\endgroup}
\newcommand{\blok}[1]{\needspace{5\baselineskip}\subsection*{#1}}
\newcommand{\vyrok}[1]{\uv{#1}}

\nastavPripravu{1. ročník}{Algebraické výrazy}{50}{Rozklad kvadratických trojčlenů}
\begin{document}
\poznamka{45minutové jádro: 5 min opakování vytýkání a algebraických vzorců,
10 min odvození postupu pro trojčlen, 10 min první dva řešené příklady,
15 min samostatná práce a její kontrola, 5 min shrnutí.
Záporné koeficienty a kombinaci s vytýkáním zařaď podle tempa.
Rezerva slouží k vysvětlení, proč nelze každý trojčlen rozložit na reálné lineární činitele.
Domácí práce přibližně 15--20 minut. Součet a rozdíl krychlí byly již zavedeny v hodině 46,
zde je znovu nevykládáme.}
\blok{Od roznásobení k rozkladu}
V minulých hodinách jsme mnohočleny rozkládali vytýkáním a pomocí algebraických
vzorců. Dnes budeme hledat součin dvou lineárních dvojčlenů i tam, kde na první
pohled není společný činitel ani známý vzorec. O výsledku rozhoduje kontrola
zpětným roznásobením.
\begin{prikladReseniBox}[title={Úvodní příklad: čti roznásobení obráceně}]
Roznásob $(x+2)(x+3)$. Která dvě čísla určují koeficienty u $x$ a konstantního členu?
\tcblower
\[(x+2)(x+3)=x^2+3x+2x+6=x^2+5x+6.\]
Čísla $2$ a $3$ mají součet $5$ a součin $6$. Proto můžeme také psát
$x^2+5x+6=(x+2)(x+3)$.
\end{prikladReseniBox}
\blok{Trojčlen $x^2+px+q$: hledáme součet a součin}
\begin{prikladReseniBox}[title={Příklad 1: odkud se bere postup?}]
Roznásob $(x+r)(x+s)$ a urči, kdy se tento součin rovná $x^2+px+q$.
\tcblower
\[(x+r)(x+s)=x^2+(r+s)x+rs.\]
Aby se koeficienty shodovaly, musí současně platit
\[r+s=p,\qquad rs=q.\]
\end{prikladReseniBox}
\begin{definiceBox}
\textbf{Podmínka pro rozklad kvadratického trojčlenu:}
\[x^2+px+q=(x+r)(x+s),\]
právě když $r+s=p$ a $rs=q$ (pro zvolené reálné koeficienty).
\end{definiceBox}
\poznamka{Jde o hledání čísel, ne o dosud neprobranou metodu řešení kvadratických rovnic. V prvních úlohách volíme celá $r,s$. Nelze-li najít celá čísla, \emph{neznamená to obecně}, že rozklad nad reálnými čísly není možný.}
\begin{prikladReseniBox}[title={Příklad 2: znaménka rozhodují}]
Rozlož $x^2+7x+10$ a $x^2-2x-3$.
\tcblower
První: součet $7$ a součin $10$ mají čísla $5$ a $2$:
\[x^2+7x+10=(x+5)(x+2).\]
Druhý: součet $-2$ a součin $-3$ mají čísla $-3$ a $1$:
\[x^2-2x-3=(x-3)(x+1).\]
Oba výsledky ověříme roznásobením.
\end{prikladReseniBox}
\begin{zadaniReseniBox}[title={Samostatná práce: hledáme dvě čísla}]
Rozlož, pokud to lze pomocí celých čísel $r,s$:
\begin{enumerate}[label=\alph*)]
\item $x^2+9x+20$;
\item $x^2-x-6$;
\item $x^2+3x-18$;
\item $x^2-8x+12$;
\item $x^2+x+4$ (rezerva: zdůvodni, proč zde neexistuje ani reálný lineární rozklad).
\end{enumerate}
\tcblower
\begin{enumerate}[label=\alph*)]
\item $(x+4)(x+5)$, $4+5=9$, $4\cdot5=20$;
\item $(x-3)(x+2)$, $-3+2=-1$, $-3\cdot2=-6$;
\item $(x+6)(x-3)$, $6-3=3$, $6\cdot(-3)=-18$;
\item $(x-6)(x-2)$, $-6-2=-8$, $(-6)(-2)=12$;
\item Žádná reálná $r,s$ s $r+s=1$, $rs=4$ neexistují: $(r-s)^2=(r+s)^2-4rs=1-16=-15$, spor.
\end{enumerate}
\end{zadaniReseniBox}
\blok{Volitelná rezerva}
\begin{zadaniReseniBox}[title={Rozšiřující příklad: nejdřív vytkni}]
Rozlož $2x^3+10x^2+12x$.
\tcblower
$2x(x^2+5x+6)=2x(x+2)(x+3)$; součet $2+3=5$, součin $2\cdot3=6$.
\end{zadaniReseniBox}
\blok{Domácí procvičení}
\begin{zadaniReseniBox}[title={Domácí úkol: rozklad trojčlenů}]
Rozlož na součin, pokud možno pomocí celých koeficientů. Výsledek ověř roznásobením.
\begin{enumerate}[label=\alph*)]
\item $x^2+8x+15$;
\item $x^2-5x+6$;
\item $x^2+2x-15$;
\item $x^2-10x+25$;
\item $x^2-4x-12$;
\item $x^2+3x-28$;
\item $2x^3+10x^2+12x$ (nejdříve vytkni);
\item $x^2+2x+2$ (rozhodni, zda lze rozložit na reálné lineární činitele).
\end{enumerate}
\tcblower
\begin{enumerate}[label=\alph*)]
\item $(x+3)(x+5)$;
\item $(x-2)(x-3)$;
\item $(x+5)(x-3)$;
\item $(x-5)^2$;
\item $(x-6)(x+2)$;
\item $(x+7)(x-4)$;
\item $2x(x+2)(x+3)$;
\item Nelze nad $\mathbb R$ rozložit na součin dvou lineárních činitelů:
$x^2+2x+2=(x+1)^2+1>0$ pro všechna $x\in\mathbb R$.
\end{enumerate}
\end{zadaniReseniBox}
\end{document}
